\documentclass{amsart}
% For dealing with references we use the comment environment
\usepackage{verbatim}
\newenvironment{reference}{\comment}{\endcomment}
%\newenvironment{reference}{}{}
\newenvironment{slogan}{\comment}{\endcomment}
\newenvironment{history}{\comment}{\endcomment}

% For commutative diagrams we use Xy-pic
\usepackage[all]{xy}

% We use 2cell for 2-commutative diagrams.
\xyoption{2cell}
\UseAllTwocells

% We use multicol for the list of chapters between chapters
\usepackage{multicol}

% This is generally recommended for better output
\usepackage{lmodern}
\usepackage[T1]{fontenc}

% For cross-file-references

% Package for hypertext links:
\usepackage{hyperref}

% For any local file, say "hello.tex" you want to link to please

% Theorem environments.
%
\theoremstyle{plain}
\newtheorem{theorem}[subsection]{Theorem}
\newtheorem{proposition}[subsection]{Proposition}
\newtheorem{lemma}[subsection]{Lemma}

\theoremstyle{definition}
\newtheorem{definition}[subsection]{Definition}
\newtheorem{example}[subsection]{Example}
\newtheorem{exercise}[subsection]{Exercise}
\newtheorem{situation}[subsection]{Situation}

\theoremstyle{remark}
\newtheorem{remark}[subsection]{Remark}
\newtheorem{remarks}[subsection]{Remarks}

\numberwithin{equation}{subsection}

% Macros
%
\def\lim{\mathop{\mathrm{lim}}\nolimits}
\def\colim{\mathop{\mathrm{colim}}\nolimits}
\def\Spec{\mathop{\mathrm{Spec}}}
\def\Hom{\mathop{\mathrm{Hom}}\nolimits}
\def\Ext{\mathop{\mathrm{Ext}}\nolimits}
\def\SheafHom{\mathop{\mathcal{H}\!\mathit{om}}\nolimits}
\def\SheafExt{\mathop{\mathcal{E}\!\mathit{xt}}\nolimits}
\def\Sch{\mathit{Sch}}
\def\Mor{\mathop{\mathrm{Mor}}\nolimits}
\def\Ob{\mathop{\mathrm{Ob}}\nolimits}
\def\Sh{\mathop{\mathit{Sh}}\nolimits}
\def\NL{\mathop{N\!L}\nolimits}
\def\CH{\mathop{\mathrm{CH}}\nolimits}
\def\proetale{{pro\text{-}\acute{e}tale}}
\def\etale{{\acute{e}tale}}
\def\QCoh{\mathit{QCoh}}
\def\Ker{\mathop{\mathrm{Ker}}}
\def\Im{\mathop{\mathrm{Im}}}
\def\Coker{\mathop{\mathrm{Coker}}}
\def\Coim{\mathop{\mathrm{Coim}}}

% Boxtimes
%
\DeclareMathSymbol{\boxtimes}{\mathbin}{AMSa}{"02}

%
% Macros for moduli stacks/spaces
%
\def\QCohstack{\mathcal{QC}\!\mathit{oh}}
\def\Cohstack{\mathcal{C}\!\mathit{oh}}
\def\Spacesstack{\mathcal{S}\!\mathit{paces}}
\def\Quotfunctor{\mathrm{Quot}}
\def\Hilbfunctor{\mathrm{Hilb}}
\def\Curvesstack{\mathcal{C}\!\mathit{urves}}
\def\Polarizedstack{\mathcal{P}\!\mathit{olarized}}
\def\Complexesstack{\mathcal{C}\!\mathit{omplexes}}
% \Pic is the operator that assigns to X its picard group, usage \Pic(X)
% \Picardstack_{X/B} denotes the Picard stack of X over B
% \Picardfunctor_{X/B} denotes the Picard functor of X over B
\def\Pic{\mathop{\mathrm{Pic}}\nolimits}
\def\Picardstack{\mathcal{P}\!\mathit{ic}}
\def\Picardfunctor{\mathrm{Pic}}
\def\Deformationcategory{\mathcal{D}\!\mathit{ef}}

\setlength{\emergencystretch}{3em}
\begin{document}
\title[Stacks Project, Tag 07VK]{375 --- Proper flat coherent cohomology is perfect and obeys arbitrary derived base change}
\maketitle
\noindent Copyright (C) 2005--2025 Johan de Jong. Stacks Project authors. Source: \href{https://stacks.math.columbia.edu/tag/07VK}{Tag 07VK}.

\noindent Editorial difficulty note (not author text). Difficulty H2: Flat base change alone does not cover an arbitrary base algebra and does not prove perfectness. The proof replaces the Cech resolution by its bounded alternating complex, proves flatness termwise, and uses square-zero base algebras to detect uniform Tor amplitude before deducing perfectness from pseudo-coherence..

\noindent Editorial provenance note (not author text). History: excerpt from revision \texttt{a0444\allowbreak 6e57e\allowbreak c1fbc\allowbreak 252a8\allowbreak 71afc\allowbreak ec775\allowbreak 2fb28\allowbreak 07b14}, coherent.tex, lines 6171--6227. Reference presentation was adapted; original-author context and core support blocks were retained or added. The mathematical wording is unchanged. Licensed under GNU FDL 1.2 or later; no invariant sections or cover texts. See ../licenses/COPYING. Context blocks reproduce author definitions and supporting results; they are not additional counted problems.

\begin{definition}
\label{more-algebra-definition-tor-amplitude}
Let $R$ be a ring. Denote $D(R)$ its derived category.
Let $a, b \in \mathbf{Z}$.
\begin{enumerate}
\item An object $K^\bullet$ of $D(R)$ has
{\it tor-amplitude in $[a, b]$}
if $H^i(K^\bullet \otimes_R^\mathbf{L} M) = 0$ for all $R$-modules
$M$ and all $i \not \in [a, b]$.
\item An object $K^\bullet$ of $D(R)$ has {\it finite tor dimension}
if it has tor-amplitude in $[a, b]$ for some $a, b$.
\item An $R$-module $M$ has {\it tor dimension $\leq d$}
if $M[0]$ as an object of $D(R)$ has tor-amplitude in $[-d, 0]$.
\item An $R$-module $M$ has {\it finite tor dimension}
if $M[0]$ as an object of $D(R)$ has finite tor dimension.
\end{enumerate}
\end{definition}

\begin{lemma}
\label{lemma-perfect-direct-image}
Let $A$ be a Noetherian ring and set $S = \Spec(A)$. Let $f : X \to S$ be a
proper morphism of schemes. Let $\mathcal{F}$ be a coherent
$\mathcal{O}_X$-module flat over $S$. Then
\begin{enumerate}
\item $R\Gamma(X, \mathcal{F})$ is a perfect object of $D(A)$, and
\item for any ring map $A \to A'$ the base change map
$$
R\Gamma(X, \mathcal{F}) \otimes_A^{\mathbf{L}} A'
\longrightarrow
R\Gamma(X_{A'}, \mathcal{F}_{A'})
$$
is an isomorphism.
\end{enumerate}
\end{lemma}

\begin{proof}
Choose a finite affine open covering $X = \bigcup_{i = 1, \ldots, n} U_i$.
By Lemmas \href{https://stacks.math.columbia.edu/tag/01XL}{lemma separated case relative cech (Tag 01XL)} and
\href{https://stacks.math.columbia.edu/tag/01XM}{lemma base change complex (Tag 01XM)} the {\v C}ech complex
$K^\bullet = {\check C}^\bullet(\mathcal{U}, \mathcal{F})$ satisfies
$$
K^\bullet \otimes_A A' = R\Gamma(X_{A'}, \mathcal{F}_{A'})
$$
for all ring maps $A \to A'$. Let
$K_{alt}^\bullet = {\check C}_{alt}^\bullet(\mathcal{U}, \mathcal{F})$
be the alternating {\v C}ech complex. By
Cohomology, Lemma \href{https://stacks.math.columbia.edu/tag/01FM}{lemma alternating usual (Tag 01FM)}
there is a homotopy equivalence $K_{alt}^\bullet \to K^\bullet$
of $A$-modules. In particular, we have
$$
K_{alt}^\bullet \otimes_A A' = R\Gamma(X_{A'}, \mathcal{F}_{A'})
$$
as well. Since $\mathcal{F}$ is flat over $A$ we see that each $K_{alt}^n$
is flat over $A$ (see
Morphisms, Lemma \href{https://stacks.math.columbia.edu/tag/01U4}{lemma flat module characterize (Tag 01U4)}).
Since moreover $K_{alt}^\bullet$ is bounded above (this is why we switched
to the alternating {\v C}ech complex)
$K_{alt}^\bullet \otimes_A A' = K_{alt}^\bullet \otimes_A^{\mathbf{L}} A'$
by the definition of derived tensor products (see
More on Algebra, Section \href{https://stacks.math.columbia.edu/tag/06XY}{section derived tensor product (Tag 06XY)}).
By
Lemma \href{https://stacks.math.columbia.edu/tag/02O6}{lemma proper over affine cohomology finite (Tag 02O6)}
the cohomology groups $H^i(K_{alt}^\bullet)$ are finite $A$-modules.
As $K_{alt}^\bullet$ is bounded, we conclude that $K_{alt}^\bullet$
is pseudo-coherent, see
More on Algebra, Lemma \href{https://stacks.math.columbia.edu/tag/066E}{lemma Noetherian pseudo coherent (Tag 066E)}.
Given any $A$-module $M$ set $A' = A \oplus M$ where $M$ is a square zero
ideal, i.e., $(a, m) \cdot (a', m') = (aa', am' + a'm)$. By the
above we see that $K_{alt}^\bullet \otimes_A^\mathbf{L} A'$ has cohomology
in degrees $0, \ldots, n$. Hence $K_{alt}^\bullet \otimes_A^\mathbf{L} M$
has cohomology in degrees $0, \ldots, n$. Hence $K_{alt}^\bullet$ has
finite Tor dimension, see
More on Algebra, Definition \ref{more-algebra-definition-tor-amplitude}.
We win by More on Algebra, Lemma \href{https://stacks.math.columbia.edu/tag/0658}{lemma perfect (Tag 0658)}.
\end{proof}
\end{document}
